\chapter{CI/CD}

\section{GitHub Actions Overview}

OmniLaTeX uses GitHub Actions with nine workflows covering every stage of the
development lifecycle.  Each workflow is triggered by different events:

\begin{table}[htbp]
  \centering
  \begin{tabular}{@{} l l l @{}}
    \toprule
    \textbf{Workflow} & \textbf{Trigger} & \textbf{Purpose} \\
    \midrule
    \texttt{build.yml} & push, PR, manual & Full build, test, deploy \\
    \texttt{lean4-ci.yml} & push, PR, manual & Lean~4 proof verification \\
    \texttt{ctan-upload.yml} & tag push (\texttt{v*}) & CTAN package upload \\
    \texttt{ctan-release.yml} & tag push (\texttt{v*}) & GitHub Release creation \\
    \texttt{ctan.yml} & push, tag, manual & CTAN zip validation \\
    \texttt{cross-platform.yml} & push, PR, manual & Linux + Windows Docker \\
    \texttt{docker-ci.yml} & push, tag, PR & Docker image build + push \\
    \texttt{docker-digest-sync.yml} & workflow\_run & Digest sync PR \\
    \texttt{integration-matrix.yml} & weekly cron, manual & Doctype $\times$ language matrix \\
    \bottomrule
  \end{tabular}
  \caption{GitHub Actions workflow summary}
\end{table}

\section{Build Pipeline}

The primary workflow (\texttt{build.yml}) runs on every push to
\texttt{main}, \texttt{master}, or \texttt{develop}, and on pull requests
to \texttt{main}/\texttt{master}.  It consists of six jobs:

\begin{enumerate}
  \item \textbf{build} -- Pulls the pinned Docker image from
        \texttt{.env.docker}, runs
        \texttt{python3 build.py --mode prod --verbose build-all} to compile
        all 43 examples, verifies every expected PDF was produced, uploads
        root and example PDFs as artifacts, and prepares a GitHub Pages
        bundle.

  \item \textbf{performance} -- Depends on \texttt{build}.  Rebuilds all
        examples with \texttt{--timings --force}, generates a performance
        summary with regression detection via
        \texttt{.github/scripts/extract_metrics.py}, and uploads
        \texttt{build/metrics.json} as an artifact.

  \item \textbf{determinism} -- Runs independently.  Builds the thesis
        example twice with \texttt{--source-date-epoch 1700000000},
        compares page counts and file sizes with a 3\% tolerance to detect
        non-deterministic output.

  \item \textbf{test} -- Depends on \texttt{build}.  Downloads PDF
        artifacts and runs integration tests via \texttt{pytest} against
        \texttt{tests/test\_modules.py} and \texttt{tests/test\_ctan.py}.

  \item \textbf{deploy-pages} -- Depends on \texttt{build}.  Deploys the
        site bundle to GitHub Pages (only on \texttt{main} push).

  \item \textbf{changelog} -- Runs only on pull requests.  Checks that
        \texttt{CHANGELOG.md} was updated whenever \texttt{.sty} or
        \texttt{.cls} files were modified.
\end{enumerate}

\section{Lean 4 Verification}

The \texttt{lean4-ci.yml} workflow runs on pushes and pull requests to
\texttt{main}/\texttt{master}.  It uses Nix for a reproducible environment:

\begin{minted}{bash}
nix develop --command bash -c "cd specs/proofs && lake build"
\end{minted}

After building, a verification step generates a summary table for all ten
proof modules.  If any file contains \texttt{sorry}, a warning is emitted.
The ten modules verified are:

\begin{table}[htbp]
  \centering
  \begin{tabular}{@{} l l @{}}
    \toprule
    \textbf{Module} & \textbf{Domain} \\
    \midrule
    BuildModes & Build mode properties \\
    BuildSystem & Cache correctness, parallelism bounds \\
    DoctypeClassMapping & Doctype-to-class mapping (60 aliases) \\
    DocumentSettings & KOMA class partition \\
    DoctypeResolution & Alias resolution \\
    FloatInvariant & Float placement bounds \\
    FontHierarchy & Strict total order on 9 font sizes \\
    I18nCompleteness & Translation key parity (846 keys) \\
    LanguageFallback & Fallback chain stability \\
    PageGeometry & Page balance equations \\
    \bottomrule
  \end{tabular}
  \caption{Lean~4 proof modules verified in CI}
\end{table}

\section{CTAN Auto-Upload}

The \texttt{ctan-upload.yml} workflow triggers on version tags
(\texttt{v*}).  It implements a two-job pipeline backed by a five-phase
upload script (\texttt{scripts/ctan-upload.sh}):

\begin{enumerate}
  \item \textbf{Pre-flight job} -- Builds the CTAN zip, validates its
        contents (must contain \texttt{omnilatex.cls}), extracts the version
        from the tag, generates a changelog, and runs a dry-run upload to
        verify all endpoints are reachable.

  \item \textbf{Pre-flight validation} (within the upload script):
    \begin{itemize}
      \item Validate package name via
            \texttt{/upload/validatePackageName}.
      \item Validate version via
            \texttt{/upload/validatePackageVersion}.
      \item Validate CTAN path via
            \texttt{/upload/validateCtanPath} (new packages only).
      \item Check URL reachability for homepage, repository, bugs, and
            support links.
    \end{itemize}

  \item \textbf{CSRF token fetch} -- Retrieves a hidden \texttt{ticket}
        field from the CTAN upload page for form-based authentication.

  \item \textbf{Upload job} -- Rebuilds the zip, regenerates the
        changelog, and submits the package to CTAN via a multipart
        \texttt{POST} request with all required metadata fields.

  \item \textbf{Response parsing} -- Checks for success patterns
        (``has been uploaded''), failure patterns (``already exists'',
        ``invalid''), and HTTP status codes.
\end{enumerate}

Secrets required: \texttt{CTAN\_EMAIL} and \texttt{CTAN\_AUTHOR}.

\section{CTAN Release}

The \texttt{ctan-release.yml} workflow also triggers on \texttt{v*} tags.
It builds and validates the CTAN zip, then creates a GitHub Release with
auto-generated release notes via \texttt{softprops/action-gh-release@v2}.
The zip is attached to the release and uploaded as a 90-day artifact.

\section{Cross-Platform}

The \texttt{cross-platform.yml} workflow validates Docker builds on both
Linux and Windows:

\begin{itemize}
  \item \textbf{Linux job} -- Runs on \texttt{ubuntu-latest}, pulls the
        Docker image, and builds the article example in production mode.
  \item \textbf{Windows job} -- Runs on \texttt{windows-latest} with
        \texttt{continue-on-error: true} (Docker Desktop may not be
        available on all Windows runners).  Uses \texttt{linux/amd64}
        platform emulation.
  \item \textbf{Report job} -- Generates a cross-platform summary table
        showing pass/fail status for each platform.
\end{itemize}

\section{Docker CI/CD}

The \texttt{docker-ci.yml} workflow builds multi-architecture Docker images
(\texttt{linux/amd64}, \texttt{linux/arm64}) and pushes them to
\texttt{ghcr.io/wyattau/omnilatex-docker}.  Key features:

\begin{itemize}
  \item Uses BuildKit with \texttt{moby/buildkit:v0.29.0} and GHA cache
        sharing.
  \item Generates SBOM and provenance metadata for non-PR builds.
  \item Tags include \texttt{latest}, semver, and major.minor patterns.
  \item 6-hour timeout for multi-arch builds.
\end{itemize}

The companion \texttt{docker-digest-sync.yml} workflow runs automatically
after a successful Docker build.  It pulls the new \texttt{latest} image,
extracts its digest, updates \texttt{.env.docker}, and creates a pull
request with the \texttt{automation} label.  This keeps all other
workflows pointing at the latest validated image without editing workflow
files.

\section{Integration Matrix}

The \texttt{integration-matrix.yml} workflow runs weekly (Sunday 02:00~UTC)
or on manual dispatch.  It tests 13 targeted doctype $\times$ language
combinations:

\begin{table}[htbp]
  \centering
  \begin{tabular}{@{} l l @{}}
    \toprule
    \textbf{Doctype} & \textbf{Language} \\
    \midrule
    thesis & english \\
    article & english \\
    book & english \\
    presentation & english \\
    cv & english \\
    cover-letter & english \\
    dissertation & english \\
    manual & english \\
    patent & english \\
    thesis & german \\
    article & german \\
    article & french \\
    article & chinese \\
    \bottomrule
  \end{tabular}
  \caption{Integration matrix combinations}
\end{table}

Each combination generates a minimal document, compiles it with
\texttt{latexmk -lualatex}, and verifies the PDF was produced.  The
\texttt{fail-fast: false} strategy ensures all combinations run even if
some fail.

\section{Determinism}

The \texttt{determinism} job in \texttt{build.yml} verifies that the build
system produces reproducible output.  It uses \texttt{SOURCE\_DATE\_EPOCH}
to freeze timestamps:

\begin{enumerate}
  \item Build the thesis example with
        \texttt{--source-date-epoch 1700000000} and save the PDF.
  \item Run \texttt{python3 build.py clean}.
  \item Rebuild with the same epoch.
  \item Compare the two PDFs.
\end{enumerate}

Because LuaLaTeX is not fully bit-for-bit deterministic (PDF object
ordering, Lua hash table randomisation, font subsetting), the comparison
uses page counts and file sizes rather than exact SHA256 hashes.  A 3\%
file size tolerance accounts for non-determinism in font subsetting and
PDF object reordering across CI runners.
